\section{Background Theory}
\label{sec:Theoretical_Background}

%%%%%%%%%%%%%%%%%%%%%%%%%%%%%%%%%%%%%%%%%%%%%%%%%%%%%%%%%%%%%%%%%%%%%%%%%%%%

%%%%%%%%%%%%%%%%%%%%%%%%%%%%%%%%%%%%%%%%%%%%%%%%%%%%%%%%%%%%%%%%%%%%%%%%%%%%%%%
\begin{comment}
    \section{Amplifier S parameters}
Nel caso di sispositivo adattato
\begin{align*}
    S_{11}&=\dfrac{Z_{in}-Z_0}{Z_{in}+Z_0}\\
    S_{21}&=\left.\dfrac{b2}{a1}\right|_{a2=0}=-\dfrac{g_mZ_L}{2}\\
    S_{22}&=\dfrac{Z_{out}-Z_0}{Z_{out}+Z_0}
\end{align*}


\begin{circuitikz}

    \draw (0,0)--++(1,0) coordinate (p1);
        \draw (p1) to[open]++(0.25,1.25)--++ (0,-5)--++(1,-0.5)--++ (0,5)--++(-1,0.5);
    \draw (p1)--++(1,0.5)--++(0.5,-0.35) to [R]++(0,-2)--++(2,1) to[V]++(0,2)--++(-0.5,+0.35)--++(1,0.5)coordinate (p2);
    \draw (p1) to[open]++(0.25,1.25)--++ (3,1.5)--++(1,-0.5) coordinate(a)--++ (-3,-1.5);
    \draw(a)--++(0,-5)--++(-3,-1.5);

    %arrowlike shape
    \filldraw[fill=cyan!25, draw=black]
(p1) to[open] ++(0,1)
     -- ++(4,2)
     -- ++(0,-0.25)
     -- ++(-4,-2)
     -- cycle;

      %arrowlike shape
    \filldraw[fill=red!25, draw=black]
(p1) to[open] ++(0,-4)
     -- ++(4,2)
     -- ++(0,-0.25)
     -- ++(-4,-2)
     -- cycle;
     
\end{circuitikz}

\end{comment}

\subsection{Differential amplifier}
The differential amplifier single input balanced output is shown in Figure \ref{fig:Schematic_with_thevenin_equivalence}.

\begin{comment}
\begin{equation}
    V_e =V_{in}\cdot \dfrac{(R_{tail}||r_{e2})}{r_{e1}+(R_{tail}||R_{e2})}
\end{equation}
If Rtail is infinite then 
\begin{equation}
    V_e =V_{in}\cdot \dfrac{\dfrac{1}{r_{e}}}{r_{e}+\dfrac{1}{r_{e}}}=\dfrac{V_{in}}{2}
\end{equation}

\begin{circuitikz}
  \draw
  (0,0) to [sV, l=$V_s$] ++(0,-2.5)node[tlground]{};
  \draw (0,0) to [R, l=$r_{e1}$] ++(3,0) node[above]{$V_e$} coordinate(Ve)
  to [R, l=$R_{tail}$,*-] ++(0,-2.5)node[tlground]{};
  \draw (Ve) to [R, l=$r_{e2}$,-*] ++(3,0) coordinate(in2);
  \draw (0,0)to[cisourceAM,*-,invert]++(0,2);
  \draw[] (in2) to[cisourceAM,*-,invert]++(0,2);
  \draw[dashed](in2)--++(0,-2.5)node[tlground]{};

 % \draw[color=red, -{Latex[length=2mm]}, thick]     (in2) ++(-0.35,-1) node[right]{$R_s$} % punto di partenza   -- ++(0,1) arc[start angle=0, end angle=90, radius=3mm]    -- ++(-0.35,0);
\end{circuitikz}
\end{comment}

\begin{figure}[hbt!]
    \centering
        \begin{circuitikz}
          \draw
          (0,0) to [sV, l=$V_s$] ++(0,-2.5)node[tlground]{};
          \draw (0,0) --++(0.25,0)to[R,l=$R_S$]++(1,0) node[npn,anchor=B](Q){Q1};
          \draw (Q.C) to[R,l_=$R_C$]++(0,1.25)--++(0,0.25)node[above]{$VCC$}++(-0.5,0)--++(1,0);
          \draw (Q.E) --++(0,-0.25) coordinate (ReqCB) to [R, l=$R_{tail}$,*-] ++(0,-2.65)node[tlground]{};
          \draw (ReqCB)--++(2,0) to [R,l=$R_{CBeq}$]++(0,-1.25) to [V,l=\SI{0}{\volt}]++(0,-1) node[ground]{};
          
           \draw[color=black, -{Latex[length=2mm]}, thick]
            (ReqCB) ++(0.5,0.25) node[right]{CBeq}--++(0,-0.35) arc[start angle=185, end angle=270, radius=3mm]
            -- ++(0.35,0); 

    \draw[color=red,dashed, thick]
    (ReqCB) ++(1.65,0.15) arc[start angle=90, end angle=185, radius=3mm]
    -- ++(0,-2) arc[start angle=180, end angle=270, radius=3mm]--++(1.5,0)
    arc[start angle=270, end angle=360, radius=3mm]--++(0,2)
    arc[start angle=0, end angle=90, radius=3mm]--cycle;

    \begin{scope}[shift={(12,0)}]
        \draw  (0,0) to [R, l=$R_{b}$] ++(0,-2.5)node[tlground]{};
          \draw (0,0) --++(-1,0) node[npn,anchor=B,xscale=-1](Q2){\scalebox{-1}[1]{Q2}};
          \draw (Q2.C) to[R,l_=$R_C$]++(0,1.25)--++(0,0.25)node[above]{$VCC$}++(-0.5,0)--++(1,0);
          \draw (Q2.E) --++(0,-0.25) coordinate (ReqCE) to [R, l=$R_{tail}$,*-] ++(0,-2.65)node[tlground]{};
          \draw (ReqCE)--++(-2,0) to [R,l_=$R_{CEeq}$]++(0,-1.25) to [V,l_=$V_{th}$]++(0,-1) node[ground]{};
          
          \draw[color=black, -{Latex[length=2mm]}, thick]
            (ReqCE) ++(-0.5,0.25) node[left]{CEeq}--++(0,-0.35) arc[start angle=360, end angle=270, radius=3mm]
            -- ++(-0.35,0);
            
    \draw[color=red,dashed, thick]
    (ReqCE) ++(-1.65,0.15) arc[start angle=90, end angle=0, radius=3mm]
    -- ++(0,-2) arc[start angle=360, end angle=270, radius=3mm]--++(-1.5,0)
    arc[start angle=270, end angle=180, radius=3mm]--++(0,2)
    arc[start angle=180, end angle=90, radius=3mm]--cycle;
    \end{scope}
    \end{circuitikz}
    
    \caption{Schematic of the differential amplifier analysis with Thevenin equivalence}
    \label{fig:Schematic_with_thevenin_equivalence}
\end{figure}
\noindent
The Gains of the stages can be calculated by appliying Thevenin's theorem to half part of the circuit at a time \cite{Ultra Broadband Low-Power 70 GHz Active Balun in 130-nm SiGe BiCMOS} \cite{Microelectronic circuits}.
Ignoring the contributions of $R_{tail}$, the CE amplifier sees from its emitter an equivalent  source and resistance, rappresenting the other half branch. The equivalent voltage source is zero because there is not an input signal applied at the base of Q2, meanwhile the equivalent resistance matches with the input resistance of the CB:
\begin{equation}
    R_{CBeq}=\dfrac{r_{\pi}+R_b}{g_m r_{\pi}+1}
    \label{eq:R_{CBeq}}
\end{equation}
(\ref{eq:R_{CBeq}}) can also be written as $r_e\left( 1+\frac{R_b}{r_\pi}\right)$ \cite{Analysis and Design of Analog Integrated Circuits}, at the condition of $g_mr_{\pi}\gg1$.\\
The resulting CE gain is defined as:
\begin{equation}
    A_{CE0}=-\dfrac{\alpha R_C}{R_E+R_{CBeq}}=-\dfrac{\alpha R_C}{\left(r_e+\dfrac{R_S}{g_m r_{\pi}+1}\right)+\dfrac{r_{\pi}+R_b}{g_m r_{\pi}+1}}
    \label{eq:A_{CE}}
\end{equation}
\noindent
The CB can be analyzed in the same way of the CE with the difference for which this time with $V_{th}\approx V_s$. The voltage gain can be calculated as the result of two blocks connected in cascade \cite{Ultra Broadband Low-Power 70 GHz Active Balun in 130-nm SiGe BiCMOS}: the Thevenin block, which contributes with a value $A_{path}=\frac{V_s}{V_e}$, and the remaining block consisting of the open circuit voltage gain of the CB $A_{CBv0}=\frac{V_o}{V_e}$. 
The resulting gain for the CB stage is defined as:
\begin{equation}
    A_{CB0}=A_{path}\cdot A_{CBv0}= \dfrac{\alpha R_C}{\left(r_e+\dfrac{R_b}{g_m r_{\pi}+1}\right)+\dfrac{r_{\pi}+R_S}{g_m r_{\pi}+1}}
    \label{eq:A_{CB}}
\end{equation}
\noindent
(\ref{eq:A_{CE}}) and (\ref{eq:A_{CB}}) are equal in module and opposite in phase, effectively creating a balanced output signal. These results are only valid at low frequencies. in the D band, The the gain is calculated in the same but the instrinsic capacitors must be taken into account. The new impedances changes according to the topology of amplifier is taken into consideration. For the CE amplifier is know that:
\begin{align}
    Z_{\pi,CEeff}&= r_{\pi} \, || \, (sC_{be}+sC_{bc}(1+|A_{CE0}|))\\
    Z_{E,CEeff}  &= r_e+\dfrac{R_S}{g_m  Z_{\pi,CEeff}+1}\\
    Z_{C,CEeff}  &=Z_C \,|| \,sC_{bc}\left(1+\dfrac{1}{|A_{CE0}|}\right)
\end{align}
\noindent
the same results are applied for the CB amplifier but without the miller
\begin{align}
    Z_{\pi,CBeff}&= r_{\pi} \, || \, sC_{be}\\
    Z_{E,CBeff} &= r_e+\dfrac{R_b}{g_m  Z_{\pi,CBeff}+1}\\
    Z_{C,CBeff}  &= Z_C\, || \, sC_{bc}
\end{align}
\noindent
(\ref{eq:A_{CE}}) and (\ref{eq:A_{CB}}) can be updated considering that
\begin{align}
    ZCBeq&=\dfrac{Z_{\pi,CBeff}+R_b}{g_mZ_{\pi,CBeff}+1} & ZCEeq&=\dfrac{Z_{\pi,CEeff}+R_b}{g_mZ_{\pi,CEeff}+1}
\end{align}
and so the gain is:
\begin{equation}
    A_{CE}(\omega)=-\dfrac{\alpha Z_{C,CEeff}}{\left(r_e+\dfrac{R_S}{g_m Z_{\pi,CEeff}+1}\right)+\dfrac{Z_{\pi,CEeff}+R_b}{g_m Z_{\pi,CEeff}+1}}
    \label{eq:A_{CEeff}}
\end{equation}


\begin{equation}
    A_{CB}(\omega)=\dfrac{\alpha Z_{C,CBeff}}{\left(r_e+\dfrac{R_b}{g_m Z_{\pi,CBeff}+1}\right)+\dfrac{Z_{\pi,CBeff}+R_S}{g_m Z_{\pi,CBeff}+1}}
    \label{eq:A_{CBeff}}
\end{equation}
\cleardoublepage
\noindent
 The voltage gain of a network can also be expressed as a function of Z, Y, or other parameters. In particular, the voltage gain of a two-port network \cite{High-Frequency Integrated Circuits}, terminated with a load admittance (commonly $Y_L=\SI{0.02}{\siemens}$), is given by
\begin{align}
    A_v=\dfrac{-Y_{21}}{Y_{22}+Y_L}
    \label{eq:TwoPortAv}
\end{align}
For a multi-port network, and in particular for a three-port, \eqref{eq:TwoPortAv} can be extended by simultaneously loading the output ports such that $I_2=-Y_LV_2$ and $I_3=-Y_LV_3$. This leads to the complete expression for the amplitude imbalance, expressed in linear scale, of a single-input balanced-output differential amplifier:
\begin{align}
    \dfrac{A_{vCE}}{A_{vCB}}=-\dfrac{Y_{21}}{Y_{31}}\cdot\dfrac{Y_{33}+Y_L-Y_{31}Y_{23}}{Y_{22}+Y_L-Y_{21}Y_{32}}
    \label{eq:DiffGainBalanceYportsComplete}
\end{align}
In a first-order approximation, by neglecting the contributions $Y_{31}Y_{23}$ and $Y_{21}Y_{32}$, the amplitude imbalance can be rewritten as in \eqref{eq:DiffGainBalanceYportsApprox} where, in short, it depends on the product of two ratios: the ratio of the forward gains and the ratio of the output load terms.
\begin{align}
    \left| \dfrac{A_{vCE}}{A_{vCB}}=\dfrac{Y_{21}}{Y_{31}}\cdot\dfrac{Y_{33}+Y_L}{Y_{22}+Y_L}\right|
    \label{eq:DiffGainBalanceYportsApprox}
\end{align}

The simplest forms of the Y-parameters are those of the intrinsic CE and CB small-signal $\pi$-models \cite{High-Frequency Integrated Circuits}. 
\begin{comment}
    \begin{figure}
    \centering
    \begin{circuitikz}
        \draw(0,0)node[left]{E}--++(1,0)coordinate (e);
        \draw (e) to [C,l=$C_\pi$]++(0,-1.5)--++(1,0)coordinate(b) node[below]{B};
        \draw (e) to [cI,invert]++(2,0) coordinate(c)--++(1,0)node[right]{C};
        \draw (e)--++(0,0.75) to [R]++(2,0)-|(c);
        \draw (b)--++(1,0) to [C,l=$C_\mu$]++(0,1.5);
       % \draw (b)to[generic]++(0,-1.25)node[ground]{};
    \end{circuitikz}
    \caption{Caption}
    \label{fig:placeholder}
\end{figure}
\end{comment}
With further simplifications, the following expressions are obtained:
\begin{align}
  Y_{CE}&=    \begin{pmatrix}
Y_{\pi}+j\omega C_{bc} & -j\omega C_{bc} \\
g_m -j\omega C_{bc} & g_0+j\omega C_{bc}
\end{pmatrix} &
  Y_{CB}&=  \begin{pmatrix}
g_m+g_0+Y_{\pi} & -g_0 \\
-g_m -g_0 & g_0 +j\omega C_{bc}
\label{eq:Y_{CB}}
\end{pmatrix}
\end{align}
\noindent
For the purposes of this work, only the CB case is considered, and the following analysis provides an approximate description of the behaviour of its parameters.\\
The intrinsic small-signal $\pi$-model can be extended to an extrinsic model by including an external base admittance $Y_B$. The resulting elements of the Y-parameter matrix are denoted by $Y_{ji,\mathrm{eff}}$.
The equations for the CB configuration are reported in Table \ref{tab:yeff}, together with the small-signal circuits used for their derivation, from which further conclusions can be drawn. For large values of $Y_B$, such as when a large capacitor is used, $Y_{11,\mathrm{eff}}$ and $Y_{22,\mathrm{eff}}$ tend toward their maximum values, $g_0+g_m+Y_\pi$ and $g_0+Y_\mu$, respectively. These values coincide with those of the ideal CB amplifier reported in \eqref{eq:Y_{CB}}.
On the other hand, for lower values of $j\omega C_B$, $Y_B$ decreases, and the second term in both $Y_{11,\mathrm{eff}}$ and $Y_{22,\mathrm{eff}}$ decreases as well.\\
If $Y_B$ is inductive instead, the behaviour is reversed.
\mbox{}\\
\begin{comment}
\begin{align}
    Y_{21,\mathrm{eff}} = g_{m,\mathrm{eff}} = -(g_0+g_m)
\end{align}
and
\begin{align}
    Y_{21,\mathrm{eff}} = g_{m,\mathrm{eff}} = -g_0-\dfrac{Y_\mu(g_m+Y_\pi)}{Y_{\pi}+Y_{\mu}}
\end{align}
\end{comment}

%%%%%%%%%%%%%%%%%%%%%%%%%%%%%%%%%%%%%%%%%%%%%%
%%%%%%%%%%%%%%%%%%%%%%%%%%%%%%%%%%%%%%%%%%%%%%%
\noindent
\begin{tabular}{|m{0.495\textwidth}|m{0.495\textwidth}|}

\hline
% Quadrante 0: intestazione su tutta la larghezza
\multicolumn{2}{|m{0.99\textwidth}|}{
\begin{minipage}[c][0.18\textheight][c]{0.99\textwidth}
    \centering
    \vspace{0.4cm}
\begin{circuitikz}
    \draw(0,0) 
    to[V,l=$v_1$,i^>=$i_1$,invert]++(0,3)--++(1.5,0)coordinate(1);
    \draw(1) to [generic,l=$Y_{\pi}$,v<=$v_{\pi}$,*-]++(0,-1.5) coordinate (x)
    to [generic,l=$Y_{B}$,v=$v_{x}$,*-*]++(0,-1.5) node[ground]{};
    \draw (x)--++(1.25,0) to [generic,l=$Y_{\mu}$,v<=$\vphantom{V}$,i^<=$i_\mu$,-*]++(0,-1.5) coordinate(gnd1);
    \draw (1)--++(2.95,0) coordinate(y) to [cI,l_=$g_mv_{\pi}$,invert,*-*]++(0,-3);
    \draw (y)--++(1.15,0) to [R,l=$g_0$,v<=$\vphantom{V}$,i^<=$i_0$]++(0,-3) -|(gnd1);
    \draw(gnd1) to [short,i<=$i_2$]++(-1,0)-|(0,0);
\end{circuitikz}
    \vfill
\end{minipage}
}
\\
\hline
% Quadrante A
\begin{minipage}[c][0.18\textheight][c]{\linewidth}
    \centering
    \begin{align}
        Y_{11,\mathrm{eff}} = g_0+(g_m+Y_{\pi})\dfrac{Y_B+Y_{\mu}}{Y_B+Y_{\pi}+Y_{\mu}}
    \end{align}
    
\end{minipage}
&
% Quadrante B
\begin{minipage}[c][0.18\textheight][c]{\linewidth}
    \centering

\begin{align}
        Y_{21,\mathrm{eff}} = -g_0-\dfrac{g_m(Y_B+Y_{\mu})+Y_{\pi}Y_\mu}{Y_B+Y_{\pi}+Y_{\mu}}
    \end{align}
\end{minipage}
\\
\hline
% Quadrante 0: intestazione su tutta la larghezza
\multicolumn{2}{|m{0.99\textwidth}|}{
\begin{minipage}[c][0.18\textheight][c]{0.99\textwidth}
    \centering
    \vspace{0.4cm}
\begin{circuitikz}
    \draw(0,0) 
    to[V,l_=$v_2$,i^>=$i_2$,invert]++(0,3)--++(-1.5,0)coordinate(1);
    \draw(1) to [generic,l=$Y_{\mu}$,*-]++(0,-1.5) coordinate (x)
    to [generic,l=$Y_{B}$,v=$\vphantom{V}$,*-*]++(0,-1.5) node[ground]{};
    \draw (x)--++(-1.25,0) to [generic,l=$Y_{\pi}$,v>=$v_\pi$,i^<=$i_\pi$,-*]++(0,-1.5) coordinate(gnd1);
    \draw (1)--++(-2.95,0) coordinate(y) to [cI,l=$g_mv_{\pi}$,*-*]++(0,-3);
    \draw (y)--++(-1.15,0) to [R,l=$g_0$,v<=$\vphantom{V}$,i^<=$i_0$]++(0,-3) -|(gnd1);
    \draw(gnd1) to [short,i_<=$i_1$]++(1,0)-|(0,0);

\end{circuitikz}
    \vfill
\end{minipage}
}
\\
\hline
% Quadrante C
\begin{minipage}[c][0.18\textheight][c]{\linewidth}
    \centering

\begin{align}
        Y_{12,\mathrm{eff}} = -g_0-\dfrac{g_mY_{\mu}+Y_{\pi}Y_\mu}{Y_B+Y_{\pi}+Y_{\mu}}
    \end{align}
\end{minipage}
&
% Quadrante D
\begin{minipage}[c][0.18\textheight][c]{\linewidth}
    \centering

\begin{align}
        Y_{22,\mathrm{eff}} = g_0+\dfrac{Y_{\mu}(Y_B+Y_{\pi}+g_m)}{Y_B+Y_{\pi}+Y_{\mu}}
    \end{align}
\end{minipage}
\\
\hline
\end{tabular}
\captionof{table}{Extrinsic Y parameters derived from the intrinsic BJT small-signal model with the addition of $Y_B$ at the base.}
\label{tab:yeff}
\mbox{}\\

\begin{comment}
The model can be complicated with the addition of an emitter impedance but further analysis of Y parameters will be avoided, only to concentrate to the gain.
By representing the found Y parameters as a Y-model and by adding a $Y_e=1/Ze$ at the input, the resulting gain ratio would be
..
\end{comment}
\begin{comment}
\textcolor{red}{-- questa parte la metto dopo il current mirror?---\\}
Everything seen so far is valid only at low frequencies.
At higher frequencies the parasitic contributions of the BJT's and the paths, letting alone the couplings, degrade the gain and the quality of the balun.\\
Having said that, as long as the core itself is identical, the losses of one half should be the same as the losses on the other half, leaving only few components who can effectively degrade the balance quality of the differential.
Such components can be the different influence of the bjt intrinsic capacitor such as $C_{u1}$ for the CE, which is not significant in the CB as long as the base in grounded, or the ${C_{CEtail}}$ of current source which degrades the CMRR.

Talking about the parasitics of the differential itself without the tail contribuiton, a, extention on the gain can be done. the CB input impedance from \cite{Analysis and Design of Analog Integrated Circuits} can be extended as
$Z_{CEq}= \left[1+Z_b\left(\dfrac{1}{Z_e}-g_m\right)\right]Z_e$ or $\dfrac{Z_{\pi}+Z_b}{g_m Z_{\pi}+1}$ (Zu=inf, zo=inf)
\\
Voglio calcolare il cazo Zu!=0
...
\\
\textcolor{red}{-------------------------------------------------}\\
\end{comment}
\subsection{Cascode current mirror}
\eqref{eq:A_{CE}}, \eqref{eq:A_{CB}}, \eqref{eq:A_{CEeff}} and \eqref{eq:A_{CBeff}} are made in the condition of $\mathrm{R}_\mathrm{CBeq}(\mathrm{Z}_\mathrm{CBeq})$, $\mathrm{R}_\mathrm{CEeq}(\mathrm{Z}_\mathrm{CEeq}) \ll \mathrm{R}_\mathrm{tail}(\mathrm{Z}_\mathrm{tail})$, permitting the remotion of the $R_{tail}$ influence from the equations. 
The easiest way to realize such condition is to implement a very big resistance value but such choice is never implemented as it increases the dissipated power, reduces the voltage headroom and introduces parasitics in compact designs.\\
High impedance at the emitter is commonly achieved with a current mirror, among which a cascode performs even better. Furthermore, to compensate the decline of the output impedance in D band, a shunt capacitor at the base of the Q3 transitor (Figure \ref{fig:cascodeCurrentMirror}) can be applied.
\begin{figure}[hbt!]
\begin{circuitikz}
  \draw (0,0) node[npn, anchor=E] (Q4) {Q4}
        (Q4.B) to [R,-o]++(-1.1,0)node[left]{$\mathrm{V}_\mathrm{B1}$} coordinate(x)
        (Q4.E)--++(0,0)node[ground] {};
  \draw (Q4.C) node[npn, anchor=E] (Q3) {Q3};
  \draw (Q3.B)--++(-0.05,0) to [R,-o]++(-1.1,0)node[left]{$\mathrm{V}_\mathrm{B2}$} coordinate(y);
  \draw (Q3.B)++(0,0) to [C,*-]++(0,-1.1)node[tlground]{};
  
  \draw[color=red, -{Latex[length=2mm]}, thick]
    (Q3.C) ++(1,0.2) node[right]{$\mathrm{Z}_\mathrm{tail}$}% punto di partenza (a destra e leggermente sopra)
    -- ++(-0.35,0) arc[start angle=90, end angle=185, radius=3mm]
    -- ++(0,-0.45); % piccola discesa in verticale

   
\end{circuitikz}
\centering
    \caption{Schematic of the cascode current mirror without the left side biasing network}
    \label{fig:cascodeCurrentMirror}
\end{figure}
\begin{comment}
To demonstrate the previous statement, a quick analisys with Figure XX can be performed. The figure XX is the $\pi$ small signal model of the bjt, setted up to evaluate the output impedance in a low frequency- generic-input-base-impedance situation.
By carrying out the KVL at the output, at the input node and by applying KCL at the emitter, the final result without approximations in equation XX is obtained.
    \begin{align}
v_{test} &= i_{test}r_0-g_mv_{\pi}r_0+v_e; & v_{\pi} &= -v_e\dfrac{r_{\pi}}{r_{\pi}+Z_s}; & v_e &= i_{test}[R_E\,||\,(r_{pi}+Z_s)]
\end{align}
\begin{align}
    Z_{out}=\dfrac{v_{test}}{i_{test}}=r_0\left[1+ g_m \dfrac{r_{\pi}}{r_{\pi}+Z_s} [r_0\,||\,(r_{\pi}+Z_s)]\right] + [r_0\,||\,(r_{\pi}+Z_s)]
\end{align}
where the degeneration emitter is the output resistance given by Q1, which is $R_{E}=r_{01}=r_0$.\\
\end{comment}

\begin{comment}

An explanation of this choice can be given by analyzing the output impedance of the CB amplifier whose Y parameters were extrapolated and with the addition of an emitter degeneration $R_E=r_0$.
$Y_{22,eff}$ in (X) can be studied for asyntotical values of the base admittance, which in this context is a capacitor of value $j \omega C_B$.
For greater values of $Y_B$, aka $\omega C_B$, the admittance reaches the asyntotic value of $g_0+Y_u$ which represents to a bigger value in terms on impedance when confronted with a lower values of the capacitor $C_B$. 
\end{comment}
An explanation of this choice can be given by analyzing the output impedance of a CE and a CB topology.
At higher frequencies, the discrepancy between the CE and CB output impedance becomes greater because of the different weight of the intrinsic collector base capacitance defined $C_{CB}$. In fact, If the amplifier works as a CE, the collector-base capacitance can be split with the miller theorem and the resultant output will be
\begin{align}
    \mathrm{Z}_\mathrm{out,CE}\approx r_o \parallel \frac{1}{j\omega\,C_{CB}\left(1+\frac{1}{|A_{v,CE}|}\right)}
\end{align}
The Common Base amplifier has the base ideally fixed at ground ($Z_s\approx \frac{1}{sC}$) and hence, the capacitance can be taken out to the collector as a shunt element. In this situation the output impedance is given by a parallel between the impedance and the intrinsic collector-base capacitor.
\begin{align}
    \mathrm{Z}_\mathrm{out,CB}\approx r_o \parallel \frac{1}{j\omega C_{CB}}
    \label{eq:ZoutCB}
\end{align}
As a result, the $Z_{out}$ is slightly greater at frequencies above the input pole compared to the CE amplifier. More importantly, the CB configuration does not enhance the collector–base capacitor.
To add further reasons, the implementation of the shunt capacitor permits to recreate a $\frac{1}{g_m}$ virtual load for Q1, thus creating an ideal unity gain and reducing the miller capacitance effect of this component.


\subsection{Base inductor degeneration}
 
Taking the CB amplifier and replacing the shunt capacitor with an inductor is one of the topologies used to create an oscillator. In the context of this work, the oscillator theory is not going to be covered but, still, a study is necessary to understand how suitable values of L can be used to tune the impedances while maintaining the circuit of Figure \ref{fig:impedance_tuning_method} in the amplifier configuration.

\begin{figure}[hbt!]
    \centering
    \begin{circuitikz}
        \draw(0,0)node[tlground]{};
        \draw (0,0) to[L]++(0,1.5) --++(0.5,0)node[npn,anchor=B](Q){Q};
        \draw [latex-](Q.E) ++(0,-0.1)--++(0,-0.4) node[below]{$R_{in}$};
        \draw (Q.C)--++(0.75,0)node[ground]{};
    \end{circuitikz}
    \caption{Schematic of the impedance tuning method using a base inductance}
    \label{fig:impedance_tuning_method}
\end{figure}
The input impedance can be extracted from $Y_{11,eff}$ and rewritten into the from present in \cite{High-Frequency Integrated Circuits}, chapter 10. For this purpose all the simplifying assumptions indicated by the author will be maintained, starting from ignoring the emitter impedance and the output conductance $g_0$:

\begin{align}
    Y_{11,eff}&=\left(\dfrac{g_m}{Y_\pi}+1\right)\dfrac{Y_{\pi}(Y_B+Y_{\mu})}{Y_B+Y_{\pi}+Y_{\mu}}
    %
    =\left(\dfrac{g_m}{Y_\pi}+1\right)\dfrac{1}{\dfrac{Y_B+Y_{\pi}+Y_{\mu}}{Y_{\pi}(Y_B+Y_{\mu})}}\\
    &=\left(\dfrac{g_m}{Y_\pi}+1\right)\dfrac{1}{\dfrac{1}{Y_{\pi}}+\dfrac{1}{Y_B+Y_{\mu}}}\\
    %
    Z_{in,eff}&= \dfrac{1}{Y_{1,1eff}}=  \dfrac{\dfrac{1}{Y_{\pi}}+\dfrac{1}{Y_B+Y_{\mu}}}{\left(\dfrac{g_m}{Y_\pi}+1\right)}
\end{align}
By calling $Z_{\pi}=\dfrac{1}{Y_\pi}$ and $Z_{Bx}=\dfrac{1}{Y_B+Y_\mu}$ the equation in the book \cite{High-Frequency Integrated Circuits} is obtained in (\ref{eq:Zin_oscillator})
\begin{align}
    Z_{in,eff}=\dfrac{Z_{Bx}+Z_{\pi}}{1+g_mZ_{\pi}}=\dfrac{(1+j\omega C_{\pi}Z_{Bx})(g_m+j\omega C_{\pi})}{g_m^2+\omega^2C_{\pi}^2}
    \label{eq:Zin_oscillator}
\end{align}

Following the author in \cite{High-Frequency Integrated Circuits}, by assuming $(f_T/f)^2\gg 1$, so to say $g_m^2 >> \omega^2 C_{\pi}^2$, the real part of the impedance can be further simplified.
\begin{align}
    \textbf{Re}(Z_{in,eff})\approx \dfrac{1}{g_m}-\dfrac{\omega^2C_{\pi}L_b}{g_m(1-\omega^2C_{bc}L_b)}
    \label{eq:Real_Zin_oscillator}
\end{align}
The resistance (\ref{eq:Real_Zin_oscillator}) reduces with respect to L, becoming negative for 
\begin{align}
    L>\dfrac{1}{\omega^2(C_{\pi}+C_{bc})}
    \label{eq:L_limit1}
\end{align}

The evaluation of the output impedance can be achieved by the same means, with no citation provided.
    \begin{align}
        Y_{22,eff} &=\dfrac{Y_{\mu}(Y_B+Y_{\pi}+g_m)}{Y_B+Y_{\pi}+Y_{\mu}}=\dfrac{Y_{\mu}(Y_B+Y_{\pi})}{Y_B+Y_{\pi}+Y_{\mu}}\left( 1+\dfrac{g_m}{Y_B+Y_{\pi}}\right)\\
        &=\dfrac{1}{\dfrac{1}{Y_\mu}+\dfrac{1}{Y_B+Y_\pi}}\left( 1+\dfrac{g_m}{Y_B+Y_{\pi}}\right)
    \end{align}
    then
\begin{align}
    Z_{out,eff}=\dfrac{1}{Y_{22,eff}}=\dfrac{\dfrac{1}{Y_\mu}+\dfrac{1}{Y_B+Y_\pi}}{1+\dfrac{g_m}{Y_B+Y_{\pi}}}
\end{align}
By calling $Z_\mu=\dfrac{1}{Y_\mu}$ and $Z_{Bx}=\dfrac{1}{Y_B+Y_\pi}$,
\begin{align}
    Z_{out,eff}=\dfrac{Z_{bc}+Z_{Bx}}{1+g_mZ_{Bx}}
\end{align}
\textcolor{red}{if $g_m>>1/Z_{Bx}$ then:}

\begin{align}
    Re(Z_{out,eff})=\dfrac{1}{gm}-\dfrac{1-\omega^2C_{\pi}L}{\omega^2C_{\mu}L\,g_m}
    \label{eq:Real_Zout_oscillator}
\end{align}
\textcolor{red}{From now on there is an error on the inequalities.}
the approximated resistance can increase or decrease depending by the second term. In particular, the resistor increases if the numerator is negative, hence if
\begin{align}
    L<\dfrac{1}{\omega^2C_{\pi}}
\end{align}

For semplicity we could also consider a safeguard value by placing 
\begin{align}
    L<\dfrac{1}{\omega^2(C_{\pi}+C_{\mu})}<\dfrac{1}{\omega^2C_{\pi}}
    \label{eq:L_limit2}
\end{align}

\eqref{eq:L_limit2} and \eqref{eq:L_limit1} guarantee the same conditions under which the real part of $Z_{in}$ in \eqref{eq:Real_Zin_oscillator} decreases while remaining positive, and under which the real part of $Z_{out}$ \eqref{eq:Real_Zout_oscillator} increases.


\subsection{S-parameters}
As the frequency increases, the wavelength of the electromagnetic waves reduces and becomes comparable to the circuits dimension. In this situation, the lumped theory approximation, which states that the signal is constant and measurable in every point of the space in a given moment of time, is no longer valid, introducing the concept according to which the Voltage and the current moves from one point to another as an electromagnetic wave constitued by the sum of incident and reflected terms.
Given the dependency of the waves from the space and time in which the measurements are taken, and combining this with the fact that parasitic components introduced by the measuring instruments interfere with the DUT behaviour, the physical measurements of voltages and currents lead to incorrect values and are, therefore, avoided alongside the Z, Y, ABCD and other parameters.

\begin{comment}
The S parameters are introduced as a substitute of the Z, Y and other parameters because they are reliable as they do not require an open or short load terminations but a load typically 50 ohm

The S parameters are defined as the radio of the reflected and incident wave at a specific port
\end{comment}

The Scattering parameters, S-parameters for short, are introduced for RF application purposes and are defined as the ratio between the reflected and the incident wave on one port when all the others are terminated to a $Z_0$ load, commonly \SI{50}{\ohm}, to eliminate the contribution from different sources.\\
A two port network can be described as a system of two equations and can be represented in a matricial form:
\begin{align}
\begin{bmatrix}
\mathrm{S}_{11} & \mathrm{S}_{12} \\
\mathrm{S}_{21} & \mathrm{S}_{22}
\end{bmatrix}
\cdot
\begin{bmatrix}
a_1 \\
a_2
\end{bmatrix}
=
\begin{bmatrix}
b_1 \\
b_2
\end{bmatrix}
\end{align}
 where
 \begin{align}
    \mathrm{S}_{ji}&=\left. \dfrac{b_j}{a_i}\right|_{a_j=0} 
\end{align}

The S parameters are best suited for describing the network at high frequencies because they require fixed loads and therefore are more reliable in comparison to parameters that depend on a perfect short or a perfect open circuit. Despite that, with mathematical manipulation, it is still possible to jump from S parameters to a different one or vice versa.

In this work, unless otherwise stated, the amplitude imbalance and phase imbalance parameters are defined as follows:
\begin{align}
    \Delta \varepsilon(\mathrm{dB})&=\mathrm{S}_{21}(\mathrm{dB})-\mathrm{S}_{31}(\mathrm{dB})\\
    \label{eq:DiffGainBalanceSports}
    \Delta \phi(\mathrm{deg})&=unwrap(phase(\mathrm{S}_{21}))-unwrap(phase(\mathrm{S}_{31}))
\end{align}

\subsection{Stability}
Through the use of the S parameters, it is possible to define the stability behaviour of the network taken into consideration \cite{Microwave Transistor Amplifiers Analysis and Design}\cite{Pozar2012}.
A two-port network is defined as unconditionally stable if both the magnitudes of the input and output reflection coefficients are less than unity for every source and load termination over the operating frequency range.  
If this condition is not satisfied, the two-port network is defined as conditionally unstable, or equivalently conditionally stable, meaning that at least one value of the source or load termination causes oscillations and therefore must be avoided, unless the circuit is intended to operate as an oscillator.\\

The stability of a two-port network can be described by the stability factor $K$ and $\Delta$:
\begin{align}
    K=\dfrac{1-|S_{11}|^2-|S_{22}|^2-|\Delta|^2}{2|S_{12}S_{21}|}\\
    \Delta=S_{11}S_{22}-S_{12}S_{21}
\end{align}

A two-port network is said to be unconditionally stable if $K>1$ and $|\Delta|<1$.  
Equivalently, using the Edwards-Sinsky stability factors, unconditional stability is obtained when $\mu>1$ and $\mu'>1$.  
If these conditions are not satisfied, the network is considered conditionally stable.\\

The Edwards-Sinsky criterion provides a necessary and sufficient condition for unconditional stability.  
In particular, both $\mu>1$ and $\mu'>1$ must be satisfied.  
$\mu$ evaluates the stability with respect to the source termination, while $\mu'$ evaluates the stability with respect to the load termination.
\begin{align}
    \mu=\frac{1-\left|S_{11}\right|^2}{\left|S_{22}-\Delta S_{11}^*\right|+\left|S_{12}S_{21}\right|}\\
    \mu'=\frac{1-\left|S_{22}\right|^2}{\left|S_{11}-\Delta S_{22}^*\right|+\left|S_{12}S_{21}\right|}
\end{align}

\begin{comment}
    \begin{align}
    \mu=\\
    \mu'=
\end{align}
\end{comment}


\subsection{Gain}
A network's power gain can be defined with three parameters: available power gain, Operating power gain and transducer power gain. 

The available power gain is expressed as the ratio between the power available from the network, so the power delivered from the network to the load ($P_{AVN}$), and the power available from the source ($P_{AVS}$). In other words, the $G_A$ is a parameter that depends from input matching network as well as the source impedance, while it is independent from the load.

The operating power gain $G_P$ can be seen as the opposite of $G_A$ and is defined as the ratio between the power delivered to the load ($P_{L}$) and the power input to the network ($P_{in}$), which means that it is independent from the source side and is dependent from the load matching network and termination. 

The transducer gain is defined as $P_L/P_{in}$ and therefore represents, among the three, the most complete parameter as it considers both source and load impedances \cite{Revisiting the Power Gains of a Loaded Two-Port} and so it describes the entire power transfer from input to output of a one-stage device.
In simultaneous conjugate matching conditions, $G_T$, $G_A$, and $G_P$ coincide.

\subsubsection{Stability gain dependence for simultaneous conjugated networks}

The general input reflection coefficients from the source and the load termination \cite{Microwave Transistor Amplifiers Analysis and Design}\cite{Pozar2012} are 
\begin{align}
    \Gamma_{in}&=S_{11}+\dfrac{S_{12}S_{21}\Gamma_L}{1-S_{22}\Gamma_L}\\
    \Gamma_{out}&=S_{22}+\dfrac{S_{12}S_{21}\Gamma_S}{1-S_{22}\Gamma_S}
\end{align}
When the two-port network is unilateral, so for $S_{12}=0$, the reflection parameters simplify , becoming 
$\Gamma_{in}=S_{11}$ and $\Gamma_{out}=S_{22}$.\\
In this condition $G_T$ takes the name of unilateral transducer power gain $G_{UT}$. Its maximum value, in \eqref{eq:GutMax}, is achieved by applying simultaneous conjugate matching at the input and output, so as to maximize $G_s$ and $G_L$. 
\begin{equation}
G_{UTmax}=
\underbrace{\vphantom{\dfrac{1}{1-|S_{11}|^2}}\dfrac{1}{1-|S_{11}|^2}}_{G_{Smax}}
\cdot
\underbrace{\vphantom{\dfrac{1}{1-|S_{11}|^2}}|S_{21}|^2}_{G_0}
\cdot
\underbrace{\vphantom{\dfrac{1}{1-|S_{11}|^2}}\dfrac{1}{1-|S_{22}|^2}}_{G_{Lmax}}
\label{eq:GutMax}
\end{equation}

In a realistic case study, a two-port network is not perfectly unilateral, but it can be treated as such if $S_{12}$ is sufficiently small. Furthermore, the dependence of $G_T$ on the stability factors must be taken into account.

In general, the maximum power gain is achieved when the simultaneous matching network is applied at both the source (with $\Gamma_S=\Gamma_{in}^*$ ) and load ($\Gamma_L=\Gamma_{out}^*$).\\
the reflection coefficients which permit this match can be called $\Gamma_{MS}$ and $\Gamma_{ML}$ and, from the stability point of view, the two port network is unconditionally stable if $|\Gamma_{MS}|$ and $|\Gamma_{ML}|$ are less than 1, which can be traced back to the condition $K>1$ and $|\Delta|<1$.

The Keysight ADS parameter "\textit{MaxGain}" \cite{Keysight Knowledge Center} refers to the maximum gain associated with a two-port network and, based on classical microwave amplifier theory\cite{Microwave Transistor Amplifiers Analysis and Design}\cite{Pozar2012}, it corresponds to either the maximum available gain or the maximum stable gain, depending on the stability condition of the device.
In particular, when the network is unconditionally stable, MaxGain corresponds to the MAG.
\begin{equation}
    MAG=\left.G_{T,max}\right|_{K>1}=\dfrac{|S_{21}|}{|S_{12}|}(K-\sqrt{K^2-1}) 
    \label{eq:MAG}
\end{equation}

When the network is conditionally stable, MaxGain corresponds to the MSG. 
\begin{equation}
    MSG=\left.G_{T,max}\right|_{K=1}=\dfrac{|S_{21}|}{|S_{12}|}
    \label{eq:MSG}
\end{equation}

From \eqref{eq:MAG} and \eqref{eq:MSG}, it follows that the MSG is the limiting value of the MAG for $K=1$ and, for greater values of $K$, the MAG decreases.  

\subsection{Mixed mode parameters}
The S parameters can be edited to achieve a profound analysis of differential networks. To do so, the power waves, the normalized voltages and the normalized currents of each port are rewritten by using differential mode and common mode components \cite{Microwave differential circuit design using mixed-mode S-parameters}. This permits to change the S matrix in a different form, 
taking the name of mixed mode S parameters, abbreviated as \textbf{$S_{mm}$}.

For this work, the mixed mode parameters are introduced mainly as a tool to extrapolate the CMRR of the devices, by mean of MATLAB simulation tool, S parameter ($S_{ds}$) and single input to common mode output S parameter ($S_{cs}$).

\begin{equation}
\mathbf{S}_{\mathrm{mm}} =
\begin{bmatrix}
S_{dd} & S_{dc} \\
S_{cd} & S_{cc}
\end{bmatrix}
\iff
\mathbf{S}_{\mathrm{se}} =
\begin{bmatrix}
S_{11} & S_{12} & S_{13} \\
S_{21} & S_{22} & S_{23} \\
S_{31} & S_{32} & S_{33}
\end{bmatrix}
\end{equation}

The conversion starts from a general view. The objective is to convert a three-port network, such as a balun, into a new form containing differential and common modes parameters. The conversion is performed through an algebraic transformation, obtained by multiplying the single-ended S-parameter matrix $S_{se}$ by $M$ and $M^{-1}$, where $M$ and $M^{-1}$ are the conversion matrices.

\begin{align}
\mathbf{S}_{\mathrm{mm}}
&= \mathbf{M}\,\mathbf{S}_{\mathrm{se}}\,\mathbf{M}^{-1}=
\begin{bmatrix}
1 & 0 & 0 \\
0 & \dfrac{1}{\sqrt{2}} & -\dfrac{1}{\sqrt{2}} \\
0 & \dfrac{1}{\sqrt{2}} & \dfrac{1}{\sqrt{2}}
\end{bmatrix}
\begin{bmatrix}
S_{11} & S_{12} & S_{13} \\
S_{21} & S_{22} & S_{23} \\
S_{31} & S_{32} & S_{33}
\end{bmatrix}
\begin{bmatrix}
1 & 0 & 0 \\
0 & \dfrac{1}{\sqrt{2}} & \dfrac{1}{\sqrt{2}} \\
0 & -\dfrac{1}{\sqrt{2}} & \dfrac{1}{\sqrt{2}}
\end{bmatrix}
\end{align}

Resulting in the matrix: 
\begin{equation}
    \begin{bmatrix}
S_{11} & \dfrac{1}{\sqrt{2}}(S_{12}-S_{13}) & \dfrac{1}{\sqrt{2}}(S_{12}+S_{13}) \\[1pt]
\dfrac{1}{\sqrt{2}}(S_{21}-S_{31}) & \dfrac{1}{2}(S_{22}-S_{32}-S_{23}+S_{33}) &\dfrac{1}{2}(S_{22}-S_{32}+S_{23}-S_{33}) \\[7pt]
\dfrac{1}{\sqrt{2}}(S_{12}+S_{13}) & \dfrac{1}{2}(S_{22}+S_{32}-S_{23}-S_{33}) & \dfrac{1}{2}(S_{22}+S_{32}+S_{23}+S_{33})
\end{bmatrix}
\label{eq:Mixed}
\end{equation}

which can be called as follows
\begin{equation}
    \mathbf{S}_{\mathrm{mm}} =
\begin{bmatrix}
S_{ss} & S_{sd} & S_{sc} \\
S_{ds} & S_{dd} & S_{dc} \\
S_{cs} & S_{cd} & S_{cc}
\end{bmatrix}
\end{equation}

where the most important parameters needed to CMRR are Sds and Scs
\begin{equation}
    \text{CMRR}(dB)=20\log_{10}{\dfrac{|S_{ds}|}{|S_{cs}|}}
    \label{eq:CMRR}
\end{equation}

